\documentclass[../../main.tex]{subfiles}
\graphicspath{{\subfix{../../image/}}}

\begin{document}

\section*{习题二}\label{section:ch02-习题二}

\addcontentsline{toc}{section}{习题二}

\begin{exercise}\label{exercise:ch02-01}
把字母 S、T、A、T、I、S、T、I、C、S 分别写在一张卡片上, 充分混和后重新排列, 问正好得到顺序 STATISTICS 的概率是多少？
\end{exercise}

\begin{exercise}\label{exercise:ch02-02}
若 $M$ 件产品中包含 $m$ 件废品, 今在其中任取两件, 求：（1）取出的两件中至少有一件是废品的概率; （2）已知取出的两件中有一件是废品的条件下, 另一件也是废品的条件概率; （3）已知两件中有一件不是废品的条件下, 另一件是废品的条件概率.
\end{exercise}

\begin{exercise}\label{exercise:ch02-03}
甲袋中有 $a$ 只白球, $b$ 只黑球, 乙袋中有 $\alpha$ 只白球, $\beta$ 只黑球, 某人从甲袋中任取两球投入乙袋, 然后在乙袋中任取两球, 问最后取出的两球全为白球的概率是多少？
\end{exercise}

\begin{exercise}\label{exercise:ch02-04}
设一个家庭中有 $n$ 个小孩的概率为
\begin{align*}
p_n=
\begin{cases}
\alpha p^n, & n\geqslant 1\\
1-\dfrac{\alpha p}{1-p}, & n=0
\end{cases}
\end{align*}
这里 $0<p<1$, $0<\alpha<\dfrac{1-p}{p}$, 若认为生一个小孩为男孩或女孩是等可能的, 求证一个家庭有 $k$（$k\geqslant 1$）个男孩的概率为 $2\alpha p^k/(2-p)^{k+1}$.
\end{exercise}

\begin{exercise}\label{exercise:ch02-05}
在上题假定下：（1）已知家庭中至少有一个男孩, 求此家庭至少有两个男孩的概率; （2）已知家庭中没有女孩, 求正好有一个男孩的概率.
\end{exercise}

\begin{exercise}\label{exercise:ch02-06}
已知产品中 $96\%$ 是合格的, 现有一种简化的检查方法, 它把真正的合格品确认为合格品的概率为 $0.98$, 而误认废品为合格品的概率为 $0.05$, 求以简化法检查下为合格品的一个产品确实是合格的概率.
\end{exercise}

\begin{exercise}\label{exercise:ch02-07}
炮战中, 若在距目标 250 米, 200 米, 150 米处射击的概率分别为 $0.1$, $0.7$, $0.2$, 而在各该处射击时命中目标的概率分别为 $0.05$, $0.1$, $0.2$, 现在已知目标被击毁, 求击毁目标的炮弹是由距目标 250 米处射出的概率.
\end{exercise}

\begin{exercise}\label{exercise:ch02-08}
飞机坠落在 A、B、C 三个区域之一, 营救部门判断其概率分别为 $0.7$, $0.2$, $0.1$; 用直升机搜索这些区域, 若有残骸, 被发现的概率分别为 $0.3$, $0.4$, $0.5$, 若已用直升机搜索过 A 区域及 B 区域, 没有发现残骸, 在这种情况下, 试计算飞机坠落在 C 区域的概率.
\end{exercise}

\begin{exercise}\label{exercise:ch02-09}
选择题有 4 个答案, 只有一个是正确的. 不懂的学生从中随机选择. 假定一个学生懂与不懂的概率都是 $1/2$, 求答对的学生对该题确实懂的概率.
\end{exercise}

\begin{exercise}\label{exercise:ch02-10}
甲袋中有 3 只黑球, 7 只白球; 乙袋中有 7 只黑球, 13 只白球; 丙袋中有 12 只黑球, 8 只白球. 先以 $1:2:2$ 的概率选择甲、乙、丙中的一只袋子. 再从选中的袋子中先后摸出 2 球, 求：（1）先摸到的是黑球的概率; （2）已知后摸到的是白球, 求先摸到的是黑球的概率.
\end{exercise}

\begin{exercise}\label{exercise:ch02-11}
甲、乙两人轮流射击, 先击中目标者获胜. 设甲、乙击中目标的概率分别为 $p_1$ 及 $p_2$, 甲先射, 试求甲获胜的概率.
\end{exercise}

\begin{exercise}\label{exercise:ch02-12}
飞机有三个不同的部分遭到射击, 在第一部分被击中一弹或第二部分被击中两弹, 或第三部分被击中三弹时, 飞机才能被击落, 其命中率与每一部分的面积成正比, 设三个部分的面积的百分比为 $0.1$, $0.2$, $0.7$. 若已击中两弹, 求击落飞机的概率.
\end{exercise}

\begin{exercise}\label{exercise:ch02-13}
证明：对于事件 $A$, $B$, 关系式
\begin{align*}
P^2(AB)+P^2(\bar{A}B)+P^2(A\bar{B})+P^2(\bar{A}\bar{B})=\frac{1}{4}
\end{align*}
成立的充要条件为
\begin{align*}
P(A)=P(B)=\frac{1}{2},\quad P(AB)=\frac{1}{4}
\end{align*}
\end{exercise}

\begin{exercise}\label{exercise:ch02-14}
若 $A$ 与 $B$ 独立, 证明 $\{\varnothing, A, \bar{A}, \Omega\}$ 中任何一个事件与 $\{\varnothing, B, \bar{B}, \Omega\}$ 中任何一个事件是相互独立的.
\end{exercise}

\begin{exercise}\label{exercise:ch02-15}
若 $0<P(B)<1$, 试证：

（1）$P(A\mid B)=P(A\mid\bar{B})$;

（2）$P(A\mid B)+P(\bar{A}\mid\bar{B})=1$

均为 $A$ 与 $B$ 相互独立的充要条件.
\end{exercise}

\begin{exercise}\label{exercise:ch02-16}
（费勒）抽查一个家庭, 考察两个事件, $A$：至多有一个女孩; $B$：男女孩子都有. 假设男女的出生率都是 $1/2$. 试证：对 3 个孩子之家, $A$ 与 $B$ 独立; 而对 4 个孩子之家, $A$ 与 $B$ 不独立.
\end{exercise}

\begin{exercise}\label{exercise:ch02-17}
事件 $A$, $B$, $C$ 两两独立, $ABC=\varnothing$, $P(A)=P(B)=P(C)$, 且已知 $P(A\cup B\cup C)=\dfrac{9}{16}$, 试求 $P(A)$.
\end{exercise}

\begin{exercise}\label{exercise:ch02-18}
设 $A$, $B$, $C$ 三事件相互独立, 求证（1）$A\cup B$, $AB$, $A-B$ 皆与 $C$ 独立; （2）$\bar{A}$, $\bar{B}$, $\bar{C}$ 亦相互独立.
\end{exercise}

\begin{exercise}\label{exercise:ch02-19}
证明：事件 $A_1, A_2, \cdots, A_n$ 相互独立的充要条件是下列 $2^n$ 个等式成立：
\begin{align*}
P(\hat{A}_1\hat{A}_2\cdots\hat{A}_n)=P(\hat{A}_1)P(\hat{A}_2)\cdots P(\hat{A}_n)
\end{align*}
其中 $\hat{A}_i$ 取 $A_i$ 或 $\bar{A}_i$.
\end{exercise}

\begin{exercise}\label{exercise:ch02-20}
三个工作小组独立对某个密码进行破译, 如果他们成功的概率分别为 $0.4$, $0.5$, $0.7$, 试求该密码被成功破译的概率.
\end{exercise}

\begin{exercise}\label{exercise:ch02-21}
设 $A_1, A_2, \cdots, A_n$ 相互独立, 而 $P(A_k)=p_k$, 试求：（1）所有事件全不发生的概率; （2）诸事件中至少发生其一的概率; （3）恰好发生其一的概率.
\end{exercise}

\begin{exercise}\label{exercise:ch02-22}
当元件 $K$ 或者元件 $K_1$ 及 $K_2$ 都发生故障时电路断开, 元件 $K$ 发生故障的概率等于 $0.3$, 而元件 $K_1$, $K_2$ 发生故障的概率各为 $0.2$, 求电路断开的概率.
\end{exercise}

\begin{exercise}\label{exercise:ch02-23}
说明"重复独立试验中, 小概率事件必然发生"的确切意思.
\end{exercise}

\begin{exercise}\label{exercise:ch02-24}
甲、乙、丙三人进行某项比赛, 若三人胜每局的概率相等, 比赛规定先胜三局者为整场比赛的优胜者, 若甲胜了第一、三局, 乙胜了第二局, 问丙成为整场比赛优胜者的概率是多少？
\end{exercise}

\begin{exercise}\label{exercise:ch02-25}
设实验室器皿中产生甲类细菌与乙类细菌的机会是相同的, 若某次发现产生了 $2n$ 个细菌, 求（1）至少有一个甲类细菌的概率; （2）甲、乙两类细菌各占其半的概率.
\end{exercise}

\begin{exercise}\label{exercise:ch02-26}
掷硬币出现正面的概率为 $p$, 掷了 $n$ 次, 求下列概率：（1）至少出现一次正面; （2）至少出现两次正面.
\end{exercise}

\begin{exercise}\label{exercise:ch02-27}
甲、乙均有 $n$ 个硬币, 全部掷完后分别计算掷出的正面数, 试求两人掷出的正面数相等的概率.
\end{exercise}

\begin{exercise}\label{exercise:ch02-28}
在 Bernoulli 试验中, 事件 $A$ 出现的概率为 $p$, 求在 $n$ 次独立试验中事件 $A$ 出现奇数次的概率.
\end{exercise}

\begin{exercise}\label{exercise:ch02-29}
在 Bernoulli 试验中, 若 $A$ 出现的概率为 $p$, 试证在出现 $m$ 次 $\bar{A}$ 之前出现 $n$ 次 $A$ 的概率, 即分赌注问题中甲最终取胜的概率, 可由（2.3.13）, （2.3.14）, （2.3.15）中的任一式子表出, 即它们是相等的.
\end{exercise}

\begin{exercise}\label{exercise:ch02-30}
袋中有 10 只黑球, 10 只白球, 从中将球一只只摸出, 求在第 9 次摸球时摸得第 3 只黑球的概率.
\end{exercise}

\begin{exercise}\label{exercise:ch02-31}
设有 $N$ 个袋子, 每个袋子中装有 $a$ 只黑球, $b$ 只白球, 从第一袋中取出一球放入第二袋中, 然后从第二袋中取出一球放入第三袋中, 如此下去, 问从最后一个袋中取出一球而为黑球的概率是多少？
\end{exercise}

\begin{exercise}\label{exercise:ch02-32}
甲袋中有 $N-1$ 只白球和 1 只黑球, 乙袋中有 $N$ 只白球, 每次从甲、乙两袋中分别取出一只球并交换放入另一袋中去, 这样经过了 $n$ 次, 问黑球出现在甲袋中的概率是多少, 并讨论 $n\to\infty$ 时的情况.
\end{exercise}

\begin{exercise}\label{exercise:ch02-33}
投硬币 $n$ 回, 第一回出正面的概率为 $c$, 第二回后每次出现与前一次相同表面的概率为 $p$, 求第 $n$ 回时出正面的概率, 并讨论当 $n\to\infty$ 时的情况.
\end{exercise}

\begin{exercise}\label{exercise:ch02-34}
甲、乙两袋各装一只白球一只黑球, 从两袋中各取出一球相交换放入另一袋中, 这样进行了若干次. 以 $p_n$, $q_n$, $r_n$ 分别记在第 $n$ 次交换后甲袋中将包含两只白球、一只白球一只黑球、两只黑球的概率. 试导出 $p_{n+1}$, $q_{n+1}$, $r_{n+1}$ 用 $p_n$, $q_n$, $r_n$ 表出的关系式, 利用它们求 $p_{n+1}$, $q_{n+1}$, $r_{n+1}$ 的表达式, 并讨论当 $n\to\infty$ 时的情况.
\end{exercise}

\begin{exercise}\label{exercise:ch02-35}
一个工厂出产的产品中废品率为 $0.005$, 任意取来 $1000$ 件, 试计算下面概率：（1）其中至少有两件废品; （2）其中不超过 5 件废品; （3）能以 $90\%$ 的概率希望废品件数不超过多少？
\end{exercise}

\begin{exercise}\label{exercise:ch02-36}
试给出 Poisson 试验的严格表述.
\end{exercise}

\begin{exercise}\label{exercise:ch02-37}
某厂长有 7 个顾问, 假定每个顾问贡献正确意见的百分比为 $0.6$, 现为某事可行与否而个别征求各顾问意见, 并按多数人的意见作出决策, 求作出正确决策的概率.
\end{exercise}

\begin{exercise}\label{exercise:ch02-38}
一本 500 页的书, 共有 500 个错字, 每个字等可能地出现在每一页上, 试求在给定的一页上至少有 3 个错字的概率.
\end{exercise}

\begin{exercise}\label{exercise:ch02-39}
某商店中出售某种商品, 据历史记录分析, 每月销售量服从 Poisson 分布, 参数为 7, 问在月初进货时要库存多少件此种商品, 才能以 $0.999$ 的概率充分满足顾客的需要.
\end{exercise}

\begin{exercise}\label{exercise:ch02-40}
螺丝钉生产中废品率为 $0.015$, 问一盒应装多少只才能保证每盒中至少有 100 只好螺丝钉的概率不小于 $80\%$（提示：用 Poisson 逼近, 设应装 $100+k$ 只）.
\end{exercise}

\begin{exercise}\label{exercise:ch02-41}
某疫苗中所含细菌数服从 Poisson 分布, 每 1 毫升中平均含有一个细菌, 把这种疫苗放入 5 只试管中, 每试管放 2 毫升, 试求：（1）5 只试管中都有细菌的概率; （2）至少有 3 只试管中有细菌的概率.
\end{exercise}

\begin{exercise}\label{exercise:ch02-42}
实验室器皿中产生甲、乙两类细菌的机会是相等的, 且产生 $k$ 个细菌的概率为
\begin{align*}
p_k=\frac{\lambda^k}{k!}e^{-\lambda},\quad k=0,1,2,\cdots
\end{align*}
试求：（1）产生了甲类细菌但没有乙类细菌的概率; （2）在已知产生了细菌而且没有甲类细菌的条件下, 有 2 个乙类细菌的概率.
\end{exercise}

\begin{exercise}\label{exercise:ch02-43}
若每条蚕的产卵数服从 Poisson 分布, 参数为 $\lambda$, 而每个卵变为成虫的概率为 $p$, 且各卵是否变为成虫彼此独立, 求每蚕养活 $k$ 只小蚕的概率.
\end{exercise}

\begin{exercise}\label{exercise:ch02-44}
通过某交叉路口的汽车流可看作 Poisson 过程, 若在一分钟内没有车的概率为 $0.2$, 求在 2 分钟内有多于一车的概率.
\end{exercise}

\begin{exercise}\label{exercise:ch02-45}
若已知 $t=0$ 时, 某分子与另一分子碰撞, 又知对任何 $t\geqslant 0$ 和 $\Delta t>0$, 若不管该分子在时刻 $t$ 以前是否遭受碰撞, 在 $(t,t+\Delta t)$ 中遭到碰撞的概率等于 $\lambda\Delta t+o(\Delta t)$, 试求该分子在时刻 $\tau$ 还没有再受到碰撞的概率.
\end{exercise}

\begin{exercise}\label{exercise:ch02-46}
利用概率论的想法证明下面恒等式：
\begin{align*}
\sum_{k=0}^{N}\binom{N+k}{k}\frac{1}{2^k}=2^N
\end{align*}
\end{exercise}

\begin{exercise}\label{exercise:ch02-47}
某车间宣称自己产品的合格率超过 $99\%$, 检验人员从该车间的 10000 件产品中抽查了 100 件, 发现有两件次品, 能否据此断定该车间谎报合格率？
\end{exercise}

\begin{exercise}\label{exercise:ch02-48}
产品验收方案规定：在一批 20 件产品中, 抽取其中 4 件, 若发现 1 件或 0 件次品, 则接受此批产品. 如果一批 20 件产品中含有 5 件次品, 若依上述方案验收, 试求这批产品被接受的概率.
\end{exercise}

\begin{exercise}\label{exercise:ch02-49}
系统中每个元件正常工作的概率为 $p$, 有半数元件正常则系统可工作, 对什么 $p$ 值, $2k+1$ 个元件的系统比 $2k-1$ 个元件的系统好？
\end{exercise}

\begin{exercise}\label{exercise:ch02-50}
通过构造适当的概率模型证明：从正整数中随机地选取两数, 此两数互素的概率等于 $\dfrac{6}{\pi^2}$.
\end{exercise}

\end{document}